\documentclass[11pt]{article}
\usepackage{amsmath}
\usepackage{amssymb}
\usepackage{booktabs}
\usepackage[margin=1in]{geometry}
\usepackage{hyperref}

\newcommand{\Mp}{M_{\mathrm{P}}}
\newcommand{\Om}{\Omega_{m,0}}
\newcommand{\Or}{\Omega_{r,0}}
\newcommand{\OL}{\Omega_{\Lambda,0}}
\newcommand{\ups}{\upsilon_2}

\title{A scaling symmetry of the TorC background equations:\\
why $\ups$ is not a parameter, and why only ratios of the $\Omega$'s are measurable}
\author{}
\date{\today}

\begin{document}
\maketitle

\begin{abstract}
Every term of the TorC background Friedmann equation is homogeneous of degree
two in time derivatives. This implies an exact one-parameter scaling symmetry
under which the three density parameters are multiplied by a common factor,
$H$ and $\dot\varpi$ by its square root, and $\varpi(a)$ is unchanged. The
observable $E(z)=H/H(a{=}1)$ is invariant. Three apparently distinct facts are
consequences of this single symmetry: (i) $\ups$ does not affect background
distances; (ii) $\Om$, $\Or$ and $\OL$ are measurable only through their
ratios; (iii) the posterior in the $(\Om,\OL)$ plane is an elongated wedge
pointing at the origin. All statements are verified numerically against the
solver to $\sim10^{-9}$.
\end{abstract}

\section{Setup}

The Lagrangian is
\begin{align}
\mathcal{L} &= \frac{3\Mp^2\ups}{4\varsigma}\left[\tfrac12(4\varsigma-\varpi^2)R
 - g^{\mu\nu}\partial_\mu\varpi\,\partial_\nu\varpi
 + \varsigma\phi^2 - \varpi^2\phi^2
 + \tfrac43\sqrt{J_\mu J^\mu}\right] - \Mp^2\Lambda + \mathcal{L}_{\mathrm M},
\label{eq:lagrangian}\\
J_\mu &\equiv \varpi^3\partial_\mu(\phi/\varpi) - \partial_\mu\phi .
\end{align}
Note for later that $\ups$ occurs \emph{once}, in the overall prefactor, and
nowhere inside the bracket, in $J_\mu$, or in $\mathcal{L}_{\mathrm M}$.

The background Friedmann equation is
\begin{equation}
H^2 = \frac{1}{\varpi^2-4\varsigma}\left[
 -\frac{4H_0^2\varsigma}{\ups}S(a)
 + 2H\varpi\dot\varpi
 + \dot\varpi^2
 + \frac{4\big(H\varsigma(\varpi^2-1) + \varsigma\varpi\dot\varpi\big)^2}
        {\varsigma-\varpi^2}\right],
\label{eq:friedmann}
\end{equation}
where we abbreviate the matter source
\begin{equation}
S(a) \equiv \Or a^{-4} + \Om a^{-3} + \OL .
\label{eq:source}
\end{equation}
The three density parameters appear in \eqref{eq:friedmann} only inside
$S(a)$, and $S(a)$ appears only once.

\section{The observation: every term is quadratic in rates}

Call $H$ and $\dot\varpi$ \emph{rates} (one time derivative each); $H_0$ is
also a rate. Then in \eqref{eq:friedmann}:

\begin{center}
\begin{tabular}{ll}
\toprule
term & rates \\
\midrule
$H^2$ & 2 \\
$H_0^2 S(a)/\ups$ & 2 \\
$2H\varpi\dot\varpi$ & 2 \\
$\dot\varpi^2$ & 2 \\
$4\big(H\varsigma(\varpi^2-1)+\varsigma\varpi\dot\varpi\big)^2/(\varsigma-\varpi^2)$ & 2 \\
\bottomrule
\end{tabular}
\end{center}

Every term carries exactly two. The quantities $\varpi$, $\varsigma$ and $a$
appear at many powers but carry no time derivative and never spoil the count.
Equation \eqref{eq:friedmann} is therefore homogeneous of degree $2$ in rates.

\section{The symmetry}

\begin{quote}
\textbf{Proposition.} Let $\kappa>0$ be an arbitrary constant --- a free
parameter labelling the transformation, with no physical meaning ($\kappa=1$ is
the identity). It is \emph{not} the speed of light. Then the transformation
\begin{align}
&\Or\to \kappa^2\Or,\qquad \Om\to \kappa^2\Om,\qquad \OL\to \kappa^2\OL, \notag\\
&H(a)\to \kappa H(a),\qquad \dot\varpi\to \kappa\dot\varpi,\qquad
\varpi(a)\to\varpi(a),\qquad \tau\to\tau/\kappa
\label{eq:symmetry}
\end{align}
leaves \eqref{eq:friedmann} invariant.
\end{quote}

\noindent\emph{Proof.} Under \eqref{eq:symmetry}, $S(a)\to \kappa^2S(a)$ since every
term of \eqref{eq:source} is linear in the $\Omega$'s. Then term by term:
\begin{align*}
H^2 &\to \kappa^2H^2, \\
\frac{4H_0^2\varsigma}{\ups}S &\to \kappa^2\,\frac{4H_0^2\varsigma}{\ups}S, \\
2H\varpi\dot\varpi &\to (\kappa)(\kappa)\,2H\varpi\dot\varpi = \kappa^2\,2H\varpi\dot\varpi, \\
\dot\varpi^2 &\to \kappa^2\dot\varpi^2, \\
H\varsigma(\varpi^2-1)+\varsigma\varpi\dot\varpi
 &\to \kappa\big[H\varsigma(\varpi^2-1)+\varsigma\varpi\dot\varpi\big]
 \;\Longrightarrow\; (\cdot)^2 \to \kappa^2(\cdot)^2 ,
\end{align*}
while the prefactor $1/(\varpi^2-4\varsigma)$ and the denominator
$\varsigma-\varpi^2$ are unchanged because $\varpi(a)$ is unchanged. Both sides
of \eqref{eq:friedmann} acquire exactly $\kappa^2$. $\square$

\subsection{Corollary: the observable is blind to it}

The quantity compared with data is
\begin{equation}
E(z) \equiv \frac{H(z)}{H(a{=}1)}
 \;\longrightarrow\; \frac{\kappa\,H(z)}{\kappa\,H(a{=}1)} = E(z).
\end{equation}
The common factor $\kappa$ cancels between numerator and denominator. The symmetry
multiplies $H$ by the \emph{same} factor at every epoch, and $E$ is constructed
precisely to divide out the overall rate.

\subsection{Three equivalent statements of the same symmetry}

Because the $\Omega$'s enter \eqref{eq:friedmann} only in the combination
$H_0^2\Omega/\ups$, the following are the same transformation:
\begin{equation}
\Omega_i \to \kappa^2\Omega_i
\quad\Longleftrightarrow\quad
\ups \to \ups/\kappa^2
\quad\Longleftrightarrow\quad
\tau \to \tau/\kappa .
\label{eq:three}
\end{equation}
The last form is the statement that the symmetry is a reparametrisation of
time: a rescaling of the clock, not of the physics.

\section{Consequence 1: $\ups$ is not a background parameter}

By \eqref{eq:three}, changing $\ups$ is a pure time rescaling. Writing
$\tau=\sqrt{\ups}\,\sigma$ removes $\ups$ from the equations of motion and from
the early-time series initial conditions entirely, so $a(\sigma)$ and
$\varpi(\sigma)$ are the same functions for every $\ups$. Since
$H\propto a^{-2}\,\mathrm da/\mathrm d\tau$ and
$\mathrm da/\mathrm d\tau = a'(\sigma)/\sqrt{\ups}$, we obtain
\begin{equation}
H(a;\ups) = \frac{H(a;\ups{=}1)}{\sqrt{\ups}}
\end{equation}
at \emph{every} epoch, hence $E(z)$ is independent of $\ups$.

This is visible in \eqref{eq:lagrangian}: $\ups$ appears only in the overall
prefactor $3\Mp^2\ups/4\varsigma$, so it is absorbed by the renaming
$\widetilde\Mp^2\equiv\Mp^2\ups$, $\widetilde\Lambda\equiv\Lambda/\ups$:
\begin{equation}
\mathcal{L} = \frac{3\widetilde\Mp^2}{4\varsigma}\big[\;\cdots\;\big]
 - \widetilde\Mp^2\widetilde\Lambda + \mathcal{L}_{\mathrm M}.
\end{equation}
By contrast $\varsigma$ appears in the prefactor \emph{and} in
$(4\varsigma-\varpi^2)R$ \emph{and} in $\varsigma\phi^2$; three independent
occurrences, so no single renaming removes it. That is precisely why
$\varsigma$ is a physical parameter and $\ups$ is not.

\paragraph{Practical consequence.} $\ups$ and $H_0$ enter only as $H_0^2/\ups$
and are therefore exactly degenerate. They must never be sampled together. In
the pipeline $\ups$ is fixed (\texttt{MUPS2\_FIXED}), which is necessary rather
than merely convenient. $\ups$ becomes physical at perturbation level, where a
comoving wavenumber $k$ introduces a scale that does not co-rescale with $\tau$.

\section{Consequence 2: only ratios of the $\Omega$'s are measurable}

The symmetry \eqref{eq:symmetry} multiplies all three $\Omega$'s by a common
factor. Any function of $(\Or,\Om,\OL)$ that survives in $E(z)$ must therefore
be invariant under that common rescaling, i.e.\ must be a function of the
ratios alone. Explicitly, factoring \eqref{eq:source},
\begin{equation}
S(a) = \Or\left[a^{-4} + \frac{\Om}{\Or}a^{-3} + \frac{\OL}{\Or}\right],
\end{equation}
the prefactor $\Or$ sets the \emph{height} of the source and is absorbed by
\eqref{eq:three}; the bracket sets its \emph{shape} and depends on two ratios
only. The shape is fixed by the two crossover epochs
\begin{equation}
\Or a^{-4} = \Om a^{-3} \;\Longrightarrow\; 1+z_{\mathrm{eq}} = \frac{\Om}{\Or},
\qquad
\Om a^{-3} = \OL \;\Longrightarrow\; 1+z_\Lambda = \left(\frac{\OL}{\Om}\right)^{1/3},
\end{equation}
both manifestly ratios. Three input numbers thus carry two physical
(shape) numbers plus one unobservable (height).

\section{Consequence 3: the $(\Om,\OL)$ wedge}

In the fit $\Or$ is held fixed, so the exact symmetry cannot be realised; only
the two-parameter motion
\begin{equation}
(\Om,\OL)\to(\kappa^2\Om,\,\kappa^2\OL)
\end{equation}
is available. This differs from the exact symmetry only by the piece that was
not rescaled,
\begin{equation}
\Delta S = (\kappa^2-1)\,\Or a^{-4},
\end{equation}
which is the \emph{sole} source of constraint along that direction. Over the
redshift range covered by DESI DR2 BAO, $\Or a^{-4}$ is below $0.1\%$ of $S$
(Table~\ref{tab:rad}), so the symmetry is only weakly broken and the posterior
stretches into a wedge pointing at the origin, with slope $\OL/\Om$ and angular
width set by the uncertainty on that ratio.

\section{Comparison with $\Lambda$CDM}

The same symmetry exists in $\Lambda$CDM. From
$H^2 = H_0^2(\Om a^{-3}+\Or a^{-4}+\OL)$, scaling all three $\Omega$'s by $\kappa^2$
gives $H\to cH$ and leaves $E=H/H_0$ unchanged. $\Lambda$CDM escapes the
degeneracy only because flatness,
\begin{equation}
\Om + \Or + \OL = 1,
\label{eq:flat}
\end{equation}
is a relation among the $\Omega$'s alone and forbids $\kappa\neq1$. Equation
\eqref{eq:flat} is simultaneously the Friedmann constraint at $a=1$ and the
normalisation $E(0)=1$; it supplies a canonical denominator, so that
$\Om = \Om/(\Om{+}\Or{+}\OL)$ identically and $\Om$ may be read as a fraction.

TorC possesses a constraint at $a=1$ as well, but it involves $\varpi_0$ and
$\dot\varpi_0$ and does not reduce to a relation among the $\Omega$'s. No
canonical denominator exists, so the normalisation must be fixed by convention
--- in the pipeline, by hardcoding $\Or=8.24\times10^{-5}$. Consequently $\Om$
alone is a ratio with respect to an arbitrary reference, whereas
\begin{equation}
\frac{\Om}{\Om+\Or+\OL}
\end{equation}
is invariant under \eqref{eq:symmetry} and is the correct analogue of the
$\Lambda$CDM $\Om$.

Both theories measure only ratios. The difference is that $\Lambda$CDM has a
natural denominator and TorC does not.

\section{Numerical verification}

All checks use the pipeline solver (Dopri8, \texttt{rtol}=\texttt{atol}=$10^{-8}$
unless stated); quoted residuals are at solver tolerance.

\begin{table}[h]\centering
\caption{The symmetry acting at the posterior-mean point, $\kappa^2=4$. The source
is multiplied by $\kappa^2$ and $H$ by $c$ at \emph{every} epoch, so $E$ is
unchanged.}
\begin{tabular}{lrrrrrr}
\toprule
$a$ & $S$ $($\kappa{=}1$)$ & $S$ $(\kappa^2{=}4)$ & ratio & $H$ $($\kappa{=}1$)$ & $H$ $(\kappa^2{=}4)$ & ratio \\
\midrule
$10^{-6}$ & $8.2693\times10^{19}$ & $3.3077\times10^{20}$ & 4.00 & $8.7004\times10^{11}$ & $1.7401\times10^{12}$ & 2.00000 \\
$10^{-4}$ & $1.1167\times10^{12}$ & $4.4668\times10^{12}$ & 4.00 & $1.0342\times10^{8}$ & $2.0685\times10^{8}$ & 2.00000 \\
$10^{-2}$ & $3.0094\times10^{5}$ & $1.2038\times10^{6}$ & 4.00 & $6.8963\times10^{4}$ & $1.3793\times10^{5}$ & 2.00000 \\
$10^{-1}$ & $2.9453\times10^{2}$ & $1.1781\times10^{3}$ & 4.00 & $2.2076\times10^{3}$ & $4.4152\times10^{3}$ & 2.00000 \\
$0.5$ & $3.3529$ & $13.412$ & 4.00 & $2.2513\times10^{2}$ & $4.5026\times10^{2}$ & 2.00000 \\
$1.0$ & $1.3028$ & $5.2111$ & 4.00 & $1.2582\times10^{2}$ & $2.5163\times10^{2}$ & 2.00000 \\
\bottomrule
\end{tabular}
\end{table}

\noindent The resulting $E(z)$ agree to $1.0000000$ in ratio at every tabulated
$a$; across $\kappa^2\in[0.5,10]$ the maximum deviation is
$\max|E/E_0-1|\sim10^{-9}$. Holding the two ratios fixed while varying the
overall scale by a factor $10^{4}$ ($\Om$ from $0.003$ to $30$) gives
$\max|E/E_0-1| = 4.4\times10^{-9}$.

\begin{table}[h]\centering
\caption{Radiation's share of the source over the DESI DR2 BAO range. This is
the only term breaking the $(\Om,\OL)$ degeneracy.}
\label{tab:rad}
\begin{tabular}{lrr}
\toprule
$z$ & $\Or a^{-4}$ & share of $S$ \\
\midrule
0.0 & $8.240\times10^{-5}$ & $0.0063\%$ \\
0.5 & $4.172\times10^{-4}$ & $0.0209\%$ \\
1.0 & $1.318\times10^{-3}$ & $0.0393\%$ \\
2.0 & $6.674\times10^{-3}$ & $0.0748\%$ \\
2.6 & $1.384\times10^{-2}$ & $0.0943\%$ \\
\bottomrule
\end{tabular}
\end{table}

\begin{table}[h]\centering
\caption{Cost in log-likelihood of moving along the degenerate direction.
Scaling all three $\Omega$'s is exactly free; scaling only $\Om$ and $\OL$ is
almost free. For reference $1\sigma$ corresponds to $\Delta\ln\mathcal{L}\approx0.5$.}
\begin{tabular}{lrr}
\toprule
$\kappa^2$ & $\Delta\ln\mathcal{L}$, $\Om,\OL$ only & $\Delta\ln\mathcal{L}$, all three \\
\midrule
0.70 & $-0.0063$ & $2.6\times10^{-8}$ \\
1.00 & $\phantom{-}0.0000$ & $0$ \\
1.50 & $-0.0179$ & $2.4\times10^{-8}$ \\
2.00 & $-0.0437$ & $8.8\times10^{-9}$ \\
\bottomrule
\end{tabular}
\end{table}

\paragraph{$\ups$.} Rate-homogeneity applied to the second-order form of the
equations implies the identity
\begin{equation}
f\big(a,\lambda\dot a,\varpi,\lambda\dot\varpi,\;\ups/\lambda^2\big)
 = \lambda^2\, f\big(a,\dot a,\varpi,\dot\varpi,\;\ups\big),
\label{eq:homog}
\end{equation}
valid for any $\lambda$. Equation \eqref{eq:homog}
holds to $7\times10^{-12}$ and $9\times10^{-15}$ for the two vector-field
components over 2000 random points (symbolically exact, verified with
\textsc{sympy}). The series initial conditions rescale consistently: with
$\sigma=\tau/\sqrt{\ups}$, the values of $\sigma_{\mathrm{ini}}$, $a$, $\varpi$
and $\sqrt{\ups}\,\mathrm da/\mathrm d\tau$ agree across $\ups\in[0.5,3]$ to
$1.6\times10^{-16}$. Measured $H$ ratios reproduce $\ups^{-1/2}$ to
$6\times10^{-10}$, and $\max|E/E_0-1|\le1.9\times10^{-9}$.

\paragraph{$\Lambda$CDM.} As implemented with flatness, scaling $\Om$ and $\Or$
together costs $\Delta\ln\mathcal{L}=-2472$ and $-1838$ for $\kappa^2=0.5,2$. Freeing
$\OL$ and normalising by $E(0)$ restores the symmetry exactly:
$\Delta\ln\mathcal{L}=0$ for $\kappa^2$ up to $10$.

\section{Summary of measured quantities}

\begin{table}[h]\centering
\caption{DESI DR2 BAO, 1000 live points. Invariant combinations are measurable;
$\Om$ and $\OL$ individually inherit the convention $\Or=8.24\times10^{-5}$ and
are additionally limited by the prior along the wedge ($7.3\%$ of the posterior
mass lies within $0.05$ of the $\OL$ prior ceiling of $1.5$).}
\small
\begin{tabular}{lrl}
\toprule
quantity & value & status \\
\midrule
$\Om/\Or = 1+z_{\mathrm{eq}}$ & $3552 \pm 1606$ & invariant \\
$\OL/\Om$ & $3.75 \pm 1.13$ & invariant (wedge slope) \\
$z_\Lambda = (\OL/\Om)^{1/3}-1$ & $0.539 \pm 0.150$ & invariant \\
$\Om/(\Om{+}\Or{+}\OL)$ & $0.2216 \pm 0.0477$ & invariant \\
\midrule
$\Om$ & $0.2927 \pm 0.1323$ & convention-dependent \\
$\OL$ & $1.0100 \pm 0.3475$ & convention-dependent \\
\midrule
$\Lambda$CDM $\Om$ & $0.2972 \pm 0.0087$ & invariant (flatness) \\
\bottomrule
\end{tabular}
\end{table}

\noindent The invariant $\Om/(\Om{+}\Or{+}\OL) = 0.2216\pm0.0477$ is both
better determined than $\Om$ itself ($21\%$ versus $45\%$, since taking the
ratio removes the spread along the wedge) and directly comparable with the
$\Lambda$CDM value $0.2972\pm0.0087$, from which it differs by
$\approx1.6\sigma$.

\end{document}
